\section*{結言}
今回の特徴は、長時間エージェントを支えるモデルの高性能化と、ウェイト公開・実行環境の具体化が同時に進んだ点にある。NVIDIAとMetaの発表は、大規模なMoEであっても推論効率やローカル実行を強く意識している。一方、FLUX 3 Image Kleinの投稿は、画像生成と細かな編集を一体化する方向性を示した。いずれも発表者のベンチマークや実行要件を、そのまま一般的な性能保証と受け取ることはできない。

独立評価者へのアクセス拡大、数学研究者への支援、学術研究向けモデル利用制度は、AIの普及に必要な評価・研究基盤を整える動きとしてつながる。Grok Botの利用成長は投稿者自身の主張であり、当該投稿も観測窓終了後だった。速報としての鮮度を保ちつつ、日時と根拠の確かさを区別して記録することが本号の前提である。
